\section*{PROJECTIVE LIMITS OF MEASURES.}
\phantomsection
\addcontentsline{toc}{section}{PROJECTIVE LIMITS OF MEASURES.}

This is the case of a theory that was developed especially in view of the needs of Probability Theory. The problems concerning a finite sequence of random variables \(X_{1}, \ldots, X_{n}\) are solved in principle when the law \(P_{X}\) of this sequence is known: it is a positive measure of mass 1 on \(\mathbf{R}^{n}\) , such that the probability of obtaining simultaneously the inequalities \(a_{1} \leq X_{1} \leq b_{1}, \ldots, a_{n} \leq X_{n} \leq b_{n}\) is equal to \(P_{X}(C)\) where \(C\) is the closed box \([a_{1}, b_{1}] \times \ldots \times [a_{n}, b_{n}]\) of \(\mathbf{R}^{n}\) .

In practice, the measure \(P_{X}\) has discrete support or else admits a density with respect to Lebesgue measure. When it is the case of an infinite sequence \((X_{n})_{n\geq1}\) of random variables, the law \(P_{n}\) of the partial sequence \((X_{1},\ldots,X_{n})\) is known in general for every integer \(n\geq1\) ; these given laws satisfy a compatibility condition that expresses the fact that the sequence \((P_{n})_{n\geq1}\) is a projective system of measures. Until around 1920, one defined in a more or less implicit manner the probabilities of events linked to the infinite sequence by a ``natural'' passing to the limit starting with probabilities in the finite case; thus it will be assumed that the probability that a game stops is the limit, for n tending to infinity, of the probability that it stops in at most n steps. Naturally, such a theory is hardly coherent, and nothing excludes the presence of ``paradoxes'', a same probability receiving two distinct estimations according to whether it is evaluated by one or other of two procedures each as ``natural'' as the other.

Steinhaus {[}293{]} seems to have been the first to feel the necessity of considering (for the case of the toss of a coin) not only the projective system \((P_{n})_{n\geq1}\) but its limit. Shortly before, in 1919, Daniell {[}76 d{]} had proved the existence of such projective limits in general, \(^{7}\) but this result seems to have remained unknown in Europe. It is found again in 1933 by Kolmogoroff in the work {[}183{]} where this author formulates the axiomatic conception of Probability Theory. The proofs of Daniell and Kolmogoroff use a compactness argument that rests on Dini's theorem.

The Daniell-Kolmogoroff theorem leaves nothing to be desired in the case of random sequences \((X_{n})_{n\geq1}\) , but the study of random functions undertaken starting in 1935 by Kolmogoroff, Feller and Doob conceals difficulties of quite another order. Let us consider for example an interval T of R, which represents the set of instants of observations of a ``stochastic process''; the set of possible trajectories is the product space \(R^{T}\) , considered as a projective limit of partial products \(R^{H}\) , where H runs through the set of finite subsets of T; one obtains in general a projective system of measures \((\mu_{H})\) . Kolmogoroff's theorem does indeed provide a measure on \(R^{T}\) , but it is defined on a family that is notably smaller than the Borel family. \(^{8}\) A variant of Kolmogoroff's construction, which provides a measure on a topological space, is due to Kakutani {[}177{]}, and has been rediscovered many times since: one considers \(\mu_{H}\) as a measure on \(\overline{R}^{H}\) carried by \(R^{H}\),\(^{9}\) the compact space \(E=\overline{R}^{T}\) is a projective limit of the finite products \(\overline{R}^{H}\) and a measure \(\mu\) on E can be defined as a projective limit of the \(\mu_{H}\) . But this procedure has a severe inconvenience; the elements of \(\overline{R}^{T}\) do not possess any regularity property allowing the continuation of the probabilistic study of the process --- or even simply the elimination of the parasitic values \(\pm\infty\) introduced by the compactification \(\overline{R}\) of R. It can be remedied by restricting the measure \(\mu\) of \(\overline{R}^{T}\) to a particular subspace (for example \(C(T)\) in the case of brownian movement); the fundamental difficulty arises from the fact that a functional space, even of the usual type, is not necessarily \(\mu\) -measurable in \(\overline{R}^{T}\) , and even the choice of the functional space can be in question. \(^{10}\)

A decisive step was accomplished in 1956 by Prokhorov in a work {[}254{]} which exercised a determining influence on the theory of stochastic processes. By putting into a convenient axiomatic form the methods used by Wiener in the article analysed above, he establishes a general existence theorem for projective limits of measures on functional spaces.

A more restricted class of projective systems was introduced by Bochner {[}24 b{]} in 1947; it is a case of projective systems made up of real vector spaces of finite dimension and surjective linear transformations. The projective limit of such a system is identified in a natural way with the algebraic dual \(E^*\) of a real vector space \(E\) supplied with a weak topology \(\sigma(E^*, E)\) ; a corresponding projective system of measures has a limit which is a measure \(\mu\) defined for a family which is notably smaller than the Borel family of \(E^*\) . Bochner characterised completely such ``premeasures'' by their Fourier transform, which is a function on \(E\) . But this result is hardly usable in the absence of a topology on \(E\) , in which case the possibility of considering \(\mu\) as a measure on the topological dual \(E'\) of \(E\) must be examined. In an independent way, R. Fortet and E. Mourier, while seeking to generalise to random variables with values in a Banach space certain classical results of Probability Theory (the law of large numbers, the central limit theorem) brought out also the fundamental role played by the Fourier transform in these questions. But substantial progress was only made in 1956 when Gelfand {[}125 c{]} suggested that the natural framework for the Fourier transform is not that of Banach spaces or Hilbert spaces, but that of nuclear Fréchet spaces. He conjectured that every continuous function of positive type on such a space is the Fourier transform of a measure on its dual, a result established soon after by Minlos {[}222{]}. Its importance comes above all from the fact that it is applicable to distribution spaces, and that the quasi-totality of functional spaces are Borel subspaces of the space of distributions (which thus constitutes a much better setting than \(\mathbf{R}^T\) ). \(^{11}\) The theory of random distributions is an area in full expansion, and we will be content with referring the reader to the work of Gelfand and Vilenkin {[}126{]}.

The results that we have just mentioned on projective limits use the existence of topologies on the base spaces. It can be asked if there exists an analogous theory in the case of ``abstract'' measures. Von Neumann proved already in 1935 the existence of product measures in all cases, but the discovery of a counterexample by Jessen and Andersen {[}290{]} ruined the hope that every projective system of measures admits a limit. Two palliatives were discovered: in 1949, C. Ionescu-Tulcea established the existence of countable projective limits, granted the existence of suitable decompositions, \(^{12}\) a very interesting result for the study of Markov processes; besides, it became clear that the topology of the spaces does not come in except by means of the set of compact subsets. It was therefore natural to seek to axiomatise this situation inside an abstract theory, by means of the notion of a compact class of subsets of a set. This work was done in 1953 by Marczewski (who establishes by this means a theorem on abstract projective limits) and Ryll-Nardzewski (who worked on the decomposition of measures). \(^{13}\)

